\documentclass[12pt,a4paper]{article}

%% ============================================================
%% PREAMBLE
%% ============================================================
\usepackage{amsmath}
\usepackage{amssymb}
\usepackage{geometry}
\geometry{margin=2.5cm}
\usepackage{booktabs}
\usepackage{xcolor}
\usepackage{hyperref}
\usepackage{listings}
\usepackage{pgfplots}
\pgfplotsset{compat=1.18}
\usepackage[most,breakable]{tcolorbox}
\tcbuselibrary{skins, breakable}
\usepackage{enumitem}
\usepackage{fancyhdr}
\usepackage{tikz}
\usepackage{array}

%% ---- tcolorbox environments (positional title) ----
\newtcolorbox{curiositybox}[1]{colback=yellow!10, colframe=orange!80!black, fonttitle=\bfseries, title=#1, breakable}
\newtcolorbox{infobox}[1]{colback=blue!5, colframe=blue!60!black, fonttitle=\bfseries, title=#1, breakable}
\newtcolorbox{mistakebox}[1]{colback=red!5, colframe=red!60!black, fonttitle=\bfseries, title=#1, breakable}
\newtcolorbox{learnbox}[1]{colback=green!5, colframe=green!60!black, fonttitle=\bfseries, title=#1, breakable}

%% ---- lstlisting (Python) ----
\lstset{language=Python, basicstyle=\ttfamily\small, keywordstyle=\color{blue},
        commentstyle=\color{gray}, stringstyle=\color{red},
        showstringspaces=false, breaklines=true, frame=single}

%% ---- fancyhdr ----
\pagestyle{fancy}
\fancyhf{}
\lhead{Topic 03: Differentiation}
\rhead{Unit I --- Complex Variables}
\cfoot{\thepage}

%% ---- Title ----
\title{\textbf{Topic 03: Differentiation}\\[6pt]
       \large Unit I --- Complex Variable: Differentiation\\[4pt]
       \normalsize Complex Variables and Special Functions\\
       JNTUA College of Engineering (Autonomous)\\
       B.Tech Engineering}
\date{}

%% ============================================================
\begin{document}
\maketitle
\tableofcontents
\newpage

%% ============================================================
%% OPENING HOOK
%% ============================================================
\begin{curiositybox}{Hook}
In real calculus, a derivative checks one direction of change along the number line.
In complex analysis, the increment can approach zero from infinitely many directions
in the plane. So if a derivative exists, the function has passed a much tougher test.
\end{curiositybox}

\bigskip

%% ============================================================
\section{Why Complex Differentiation Is Special}
%% ============================================================

In single-variable real calculus you can approach a limit from only two directions:
from the left or from the right. If those two one-sided limits agree, the derivative
exists. The complex plane is a two-dimensional arena; the increment $\Delta z$ can
travel toward zero along any curved or straight path. A complex function
$f : \mathbb{C} \to \mathbb{C}$ must return the same limiting ratio
$\dfrac{f(z+\Delta z)-f(z)}{\Delta z}$ no matter which path $\Delta z$ follows.
This requirement is vastly more restrictive than real differentiability.

\begin{infobox}{Real vs.\ Complex Differentiability at a Glance}
\begin{center}
\begin{tabular}{@{}lll@{}}
\toprule
\textbf{Feature} & \textbf{Real calculus} & \textbf{Complex analysis} \\
\midrule
Number of approach directions & 2 (left, right) & Infinitely many \\
Differentiability $\Rightarrow$ continuity & Yes & Yes \\
Differentiability $\Rightarrow$ analyticity & Not generally & Yes (if holds on a region) \\
Differentiable $\Rightarrow$ smooth & Not always & Infinitely smooth \\
\bottomrule
\end{tabular}
\end{center}
\end{infobox}

\begin{learnbox}{Key Takeaway}
Complex differentiability is \emph{much stricter} than real differentiability.
A function that is complex-differentiable at every point of an open region gains
extraordinary smoothness properties not seen in real analysis.
\end{learnbox}

%% ============================================================
\section{Definition of Derivative}
%% ============================================================

Let $f$ be a complex-valued function defined in a neighbourhood of $z_0 \in \mathbb{C}$.

\begin{infobox}{Definition --- Complex Derivative}
The \textbf{derivative} of $f$ at $z_0$ is defined as
\[
  f'(z_0) = \lim_{\Delta z \to 0} \frac{f(z_0 + \Delta z) - f(z_0)}{\Delta z},
\]
provided this limit \emph{exists and is the same for every path} along which
$\Delta z \to 0$ in $\mathbb{C}$.

\medskip
\noindent Here $\Delta z = \Delta x + i\,\Delta y$ is a complex increment.
The limit must be independent of the manner in which $\Delta z \to 0$.
\end{infobox}

\noindent\textbf{Why path independence is the key condition.}
Write $\Delta z = |\Delta z|e^{i\alpha}$ where $\alpha$ is the angle of approach.
For the limit to exist, the difference quotient must converge to the same complex number
regardless of $\alpha$. When this holds, we say $f$ is \emph{complex differentiable}
(or \emph{holomorphic}) at $z_0$.

\begin{learnbox}{Key Takeaway}
The complex derivative is a single-limit definition, but hidden inside it is a
constraint on \emph{all} directions of approach simultaneously.
\end{learnbox}

%% ============================================================
\section{Differentiability vs.\ Derivative at a Point}
%% ============================================================

\begin{infobox}{Terminology}
\begin{itemize}[leftmargin=1.5em]
  \item \textbf{Differentiable at $z_0$}: the limit in the definition exists at the
        single point $z_0$.
  \item \textbf{Analytic (holomorphic) at $z_0$}: $f$ is differentiable at every point
        in \emph{some open neighbourhood} of $z_0$.
  \item \textbf{Analytic on a region $D$}: $f$ is analytic at every point of $D$.
  \item \textbf{Entire function}: analytic at every point of $\mathbb{C}$.
\end{itemize}
\end{infobox}

A function may be differentiable at an isolated point without being analytic there.
For example, $f(z) = |z|^2$ is differentiable only at $z = 0$ (as we shall verify
later), but fails at every other point; it is therefore \emph{not} analytic anywhere.

\medskip
The chain of implications is:
\[
  \text{Entire} \;\Rightarrow\; \text{Analytic on } D
  \;\Rightarrow\; \text{Analytic at } z_0
  \;\Rightarrow\; \text{Differentiable at } z_0
  \;\Rightarrow\; \text{Continuous at } z_0.
\]
None of these arrows reverses in general.

\begin{learnbox}{Key Takeaway}
Differentiability at a single isolated point is much weaker than analyticity.
Always state \emph{where} a function is differentiable and whether that region is open.
\end{learnbox}

%% ============================================================
\section{Standard Differentiation Rules}
%% ============================================================

\begin{infobox}{Differentiation Rules (valid wherever the functions are differentiable)}
Let $f$ and $g$ be differentiable at $z$, and let $c \in \mathbb{C}$ be a constant.

\medskip
\begin{tabular}{@{}ll@{}}
\toprule
\textbf{Rule} & \textbf{Formula} \\
\midrule
Constant & $\dfrac{d}{dz}[c] = 0$ \\[8pt]
Constant multiple & $\dfrac{d}{dz}[c\,f(z)] = c\,f'(z)$ \\[8pt]
Sum & $\dfrac{d}{dz}[f(z)+g(z)] = f'(z)+g'(z)$ \\[8pt]
Product & $\dfrac{d}{dz}[f(z)g(z)] = f'(z)g(z)+f(z)g'(z)$ \\[8pt]
Quotient & $\dfrac{d}{dz}\!\left[\dfrac{f(z)}{g(z)}\right]
           = \dfrac{f'(z)g(z)-f(z)g'(z)}{[g(z)]^2},\quad g(z)\neq 0$ \\[10pt]
Power & $\dfrac{d}{dz}[z^n] = n z^{n-1},\quad n\in\mathbb{Z}$ \\[8pt]
Chain & $\dfrac{d}{dz}[f(g(z))] = f'(g(z))\cdot g'(z)$ \\
\bottomrule
\end{tabular}
\end{infobox}

\medskip\noindent\textbf{Example --- Product rule.}
Let $f(z) = z^3$ and $g(z) = e^z$. Then:
\[
  \frac{d}{dz}[z^3 e^z] = 3z^2 e^z + z^3 e^z = e^z(3z^2 + z^3) = z^2 e^z(3+z).
\]

\medskip\noindent\textbf{Example --- Quotient rule.}
Let $h(z) = \dfrac{z^2+1}{z-1}$, $z \neq 1$. Then:
\[
  h'(z) = \frac{2z(z-1) - (z^2+1)(1)}{(z-1)^2}
         = \frac{2z^2 - 2z - z^2 - 1}{(z-1)^2}
         = \frac{z^2 - 2z - 1}{(z-1)^2}.
\]

\medskip\noindent\textbf{Example --- Chain rule.}
Let $p(z) = (z^2+3)^5$. Set $w = z^2+3$, so $p = w^5$. Then:
\[
  p'(z) = 5w^4 \cdot 2z = 10z(z^2+3)^4.
\]

\begin{learnbox}{Key Takeaway}
All familiar real-calculus differentiation rules carry over verbatim to complex
differentiable functions because the proofs use only the algebra of limits.
\end{learnbox}

%% ============================================================
\section{Differentiation of Standard Complex Functions}
%% ============================================================

\begin{infobox}{Standard Complex Derivatives}
\begin{tabular}{@{}ll@{}}
\toprule
\textbf{Function} & \textbf{Derivative} \\
\midrule
$z^n$ ($n \in \mathbb{Z}$) & $nz^{n-1}$ \\[4pt]
$e^z$ & $e^z$ \\[4pt]
$\sin z$ & $\cos z$ \\[4pt]
$\cos z$ & $-\sin z$ \\[4pt]
$\tan z$ & $\sec^2 z$ \\[4pt]
$\sinh z$ & $\cosh z$ \\[4pt]
$\cosh z$ & $\sinh z$ \\[4pt]
$\ln z$ & $\dfrac{1}{z}$ (on each branch, away from branch cut) \\
\bottomrule
\end{tabular}
\end{infobox}

\noindent\textbf{Polynomials.}
Any polynomial $P(z) = a_n z^n + a_{n-1}z^{n-1} + \cdots + a_0$ is entire;
its derivative is $P'(z) = n a_n z^{n-1} + (n-1)a_{n-1}z^{n-2} + \cdots + a_1$.

\medskip\noindent\textbf{Rational functions.}
A rational function $R(z) = P(z)/Q(z)$ is analytic everywhere except at the zeros
of $Q(z)$.

\medskip\noindent\textbf{The exponential $e^z$.}
Since $e^z = e^x(\cos y + i\sin y)$, one checks from the definition that
$(e^z)' = e^z$ and $e^z \neq 0$ for all $z$.

\medskip\noindent\textbf{$\sin z$ and $\cos z$.}
Defined via $\sin z = \dfrac{e^{iz}-e^{-iz}}{2i}$ and
$\cos z = \dfrac{e^{iz}+e^{-iz}}{2}$, both are entire functions;
differentiating these definitions gives $(\sin z)' = \cos z$ and $(\cos z)' = -\sin z$.

\medskip\noindent\textbf{Logarithm branch issue.}
$\operatorname{Log} z = \ln|z| + i\,\operatorname{Arg}(z)$ is analytic in
$\mathbb{C} \setminus (-\infty, 0]$ (the principal branch), where the branch cut
along the negative real axis must be avoided. Its derivative there is $1/z$.
Different branches differ by $2\pi i k$ for integer $k$ and are each analytic on
their respective cut planes.

\begin{learnbox}{Key Takeaway}
Polynomials and $e^z$ are entire. Rational functions fail only at poles.
$\sin z$ and $\cos z$ are entire. The logarithm requires a branch-cut choice
and is never entire.
\end{learnbox}

%% ============================================================
\section{Why Path Independence Matters}
%% ============================================================

We now sketch why the definition forces two special conditions that connect the partial
derivatives of $u = \operatorname{Re}(f)$ and $v = \operatorname{Im}(f)$.
Write $f(z) = u(x,y) + iv(x,y)$ with $z = x + iy$.

\medskip\noindent\textbf{Path 1 --- Along the real axis} ($\Delta y = 0$, $\Delta z = \Delta x$):
\[
  f'(z) = \lim_{\Delta x \to 0} \frac{[u(x+\Delta x, y)-u(x,y)] + i[v(x+\Delta x,y)-v(x,y)]}{\Delta x}
        = \frac{\partial u}{\partial x} + i\frac{\partial v}{\partial x}.
\]

\medskip\noindent\textbf{Path 2 --- Along the imaginary axis} ($\Delta x = 0$, $\Delta z = i\,\Delta y$):
\[
  f'(z) = \lim_{\Delta y \to 0}
          \frac{[u(x,y+\Delta y)-u(x,y)] + i[v(x,y+\Delta y)-v(x,y)]}{i\,\Delta y}
        = \frac{\partial v}{\partial y} - i\frac{\partial u}{\partial y}.
\]

Since both expressions must equal the \emph{same} complex number $f'(z)$, we equate
real and imaginary parts:
\[
  \frac{\partial u}{\partial x} = \frac{\partial v}{\partial y},
  \qquad
  \frac{\partial u}{\partial y} = -\frac{\partial v}{\partial x}.
\]
These are the \textbf{Cauchy--Riemann equations}, the subject of Topic~04.

\begin{infobox}{Connection to Topic 04}
Path independence of the complex derivative is \emph{equivalent} to the
Cauchy--Riemann equations holding (under suitable continuity assumptions).
Every subsequent result in complex analysis --- from analyticity to power series
to contour integration --- rests on this two-path argument.
\end{infobox}

\begin{learnbox}{Key Takeaway}
Taking the limit along two perpendicular directions and demanding equality
yields the Cauchy--Riemann equations. This is why complex differentiability
is so restrictive and so powerful.
\end{learnbox}

%% ============================================================
\section{Worked Examples}
%% ============================================================

%% --- Example 1 ---
\subsection*{Example 1: Derivative of $z^2$ from First Principles}

Let $f(z) = z^2$. Form the difference quotient:
\[
  \frac{f(z+\Delta z)-f(z)}{\Delta z}
  = \frac{(z+\Delta z)^2 - z^2}{\Delta z}
  = \frac{z^2 + 2z\,\Delta z + (\Delta z)^2 - z^2}{\Delta z}
  = \frac{2z\,\Delta z + (\Delta z)^2}{\Delta z}
  = 2z + \Delta z.
\]
Taking $\Delta z \to 0$:
\[
  f'(z) = \lim_{\Delta z\to 0}(2z + \Delta z) = 2z.
\]
The result is independent of the path because the expression $2z + \Delta z$ depends
only on $|\Delta z|$ and $z$, not on the direction of $\Delta z$.

\begin{learnbox}{Takeaway --- Example 1}
$\dfrac{d}{dz}(z^2) = 2z$, confirmed from first principles. The algebra mirrors
real calculus, but the path independence is non-trivial and crucial.
\end{learnbox}

%% --- Example 2 ---
\subsection*{Example 2: Derivative of $1/z$ from First Principles}

Let $f(z) = 1/z$, $z \neq 0$. Then:
\[
  \frac{f(z+\Delta z)-f(z)}{\Delta z}
  = \frac{\dfrac{1}{z+\Delta z} - \dfrac{1}{z}}{\Delta z}
  = \frac{z - (z+\Delta z)}{z(z+\Delta z)\,\Delta z}
  = \frac{-\Delta z}{z(z+\Delta z)\,\Delta z}
  = \frac{-1}{z(z+\Delta z)}.
\]
As $\Delta z \to 0$:
\[
  f'(z) = \frac{-1}{z \cdot z} = -\frac{1}{z^2}.
\]

\begin{learnbox}{Takeaway --- Example 2}
$\dfrac{d}{dz}\!\left(\dfrac{1}{z}\right) = -\dfrac{1}{z^2}$ for all $z \neq 0$.
This confirms the power rule $nz^{n-1}$ for $n=-1$.
\end{learnbox}

%% --- Example 3 ---
\subsection*{Example 3: $\overline{z}$ Is Nowhere Differentiable}

Let $f(z) = \overline{z} = x - iy$. Form the difference quotient with
$\Delta z = \Delta x + i\,\Delta y$:
\[
  \frac{f(z+\Delta z)-f(z)}{\Delta z}
  = \frac{\overline{z+\Delta z} - \overline{z}}{\Delta z}
  = \frac{\overline{\Delta z}}{\Delta z}
  = \frac{\Delta x - i\,\Delta y}{\Delta x + i\,\Delta y}.
\]

\textbf{Path 1} ($\Delta y = 0$, approach along real axis):
\[
  \frac{\Delta x}{\Delta x} = 1.
\]

\textbf{Path 2} ($\Delta x = 0$, approach along imaginary axis):
\[
  \frac{-i\,\Delta y}{i\,\Delta y} = -1.
\]
The two paths give different limits ($1$ and $-1$), so the limit does not exist.
Therefore $f(z) = \overline{z}$ is \emph{nowhere differentiable}.

\begin{learnbox}{Takeaway --- Example 3}
$f(z) = \overline{z}$ fails the path-independence test at every point in $\mathbb{C}$.
Real-valued intuition that ``a simple function should be differentiable'' does not
apply in complex analysis.
\end{learnbox}

%% --- Example 4 ---
\subsection*{Example 4: $|z|^2$ Is Differentiable Only at $z = 0$}

Let $f(z) = |z|^2 = x^2 + y^2$ (real-valued). Then:
\[
  \frac{f(z+\Delta z)-f(z)}{\Delta z}
  = \frac{|z+\Delta z|^2 - |z|^2}{\Delta z}.
\]
Expand $|z+\Delta z|^2 = (z+\Delta z)\overline{(z+\Delta z)}
= z\overline{z} + z\overline{\Delta z} + \overline{z}\,\Delta z + |\Delta z|^2$:
\[
  \frac{z\overline{\Delta z} + \overline{z}\,\Delta z + |\Delta z|^2}{\Delta z}
  = z \cdot \frac{\overline{\Delta z}}{\Delta z} + \overline{z} + \frac{|\Delta z|^2}{\Delta z}.
\]
As $\Delta z \to 0$, the last term $\to 0$. But $\dfrac{\overline{\Delta z}}{\Delta z}$ is
path-dependent (from Example~3, it equals $1$ along the real axis and $-1$ along the
imaginary axis) \emph{unless} $z = 0$, in which case the term $z\cdot\frac{\overline{\Delta z}}{\Delta z}$
vanishes regardless. Therefore the limit exists only at $z = 0$, where $f'(0) = \overline{0} = 0$.

\begin{learnbox}{Takeaway --- Example 4}
$|z|^2$ is differentiable \emph{only} at the origin. It is not analytic anywhere.
This shows isolated differentiability is possible without analyticity.
\end{learnbox}

%% --- Example 5 ---
\subsection*{Example 5: Derivative of $e^{2z}$ Using the Chain Rule}

Let $f(z) = e^{2z}$. Write $f(z) = e^{w}$ with $w = 2z$:
\[
  f'(z) = e^w \cdot \frac{dw}{dz} = e^{2z} \cdot 2 = 2e^{2z}.
\]

Verification via first principles: the difference quotient is
\[
  \frac{e^{2(z+\Delta z)} - e^{2z}}{\Delta z}
  = e^{2z}\cdot\frac{e^{2\,\Delta z}-1}{\Delta z}
  \to e^{2z} \cdot 2 = 2e^{2z}
\]
using $\lim_{w\to 0}\dfrac{e^w-1}{w} = 1$ (applied with $w = 2\,\Delta z \to 0$).

\begin{learnbox}{Takeaway --- Example 5}
The chain rule works identically in complex calculus. For $e^{cz}$,
the derivative is always $c\,e^{cz}$.
\end{learnbox}

%% --- Example 6 ---
\subsection*{Example 6: Testing Differentiability via Two Paths}

Let $f(z) = \dfrac{z^2}{|z|}$ for $z \neq 0$ and $f(0) = 0$. Test differentiability at $z = 0$.

Form the difference quotient at $z = 0$:
\[
  \frac{f(\Delta z) - f(0)}{\Delta z}
  = \frac{(\Delta z)^2 / |\Delta z|}{\Delta z}
  = \frac{\Delta z}{|\Delta z|}.
\]
Write $\Delta z = r e^{i\theta}$ (polar form). Then:
\[
  \frac{\Delta z}{|\Delta z|} = \frac{r e^{i\theta}}{r} = e^{i\theta} = \cos\theta + i\sin\theta.
\]
This depends on $\theta$ (the angle of approach). For example:
\begin{itemize}
  \item $\theta = 0$ (real axis): quotient $\to 1$.
  \item $\theta = \pi/2$ (imaginary axis): quotient $\to i$.
  \item $\theta = \pi/4$: quotient $\to \dfrac{1}{\sqrt{2}} + \dfrac{i}{\sqrt{2}}$.
\end{itemize}
Since the limit depends on $\theta$, $f$ is \emph{not differentiable} at the origin.

\begin{learnbox}{Takeaway --- Example 6}
When a difference quotient reduces to $e^{i\theta}$ (angle dependent), the limit
cannot exist. Checking just one or two paths can reveal non-differentiability,
but only checking \emph{all paths} (via Cauchy--Riemann) can confirm differentiability.
\end{learnbox}

%% ============================================================
\section{Excel Example (MANDATORY)}
%% ============================================================

\subsection*{Numerical Approximation of the Difference Quotient Along Two Paths}

We approximate $f'(z_0)$ for $f(z) = z^2$ at $z_0 = 1+i$ using
$\Delta z = h$ (real) and $\Delta z = ih$ (imaginary), with $h \to 0$, and compare.

\medskip
\begin{center}
\begin{tabular}{@{}cccccc@{}}
\toprule
\textbf{Col A} & \textbf{Col B} & \textbf{Col C} & \textbf{Col D}
               & \textbf{Col E} & \textbf{Col F} \\
$h$ & $\Delta z$ (real) & $\operatorname{Re}(Q_1)$ & $\operatorname{Im}(Q_1)$
    & $\operatorname{Re}(Q_2)$ & $\operatorname{Im}(Q_2)$ \\
\midrule
0.1  & real $h$  & 2.1 & 2.0 & 2.0 & 2.1 \\
0.01 & real $h$  & 2.01 & 2.0 & 2.0 & 2.01 \\
0.001 & real $h$ & 2.001 & 2.0 & 2.0 & 2.001 \\
\bottomrule
\end{tabular}
\end{center}

\medskip
\noindent\textbf{Excel column formulas} (with $z_0 = 1+i$, so $x_0=1$, $y_0=1$):

\begin{itemize}
  \item \textbf{A2}: \texttt{0.1} (step size $h$)
  \item \textbf{A3}: \texttt{=A2/10} (decreasing $h$)
  \item \textbf{Path 1 --- $\Delta z = h$ (real):}\\
        $f(z_0+h) = (1+h+i)^2 = (1+h)^2 - 1 + 2(1+h)i$\\
        \textbf{C2}: \texttt{=((1+A2)\^{}2 - 1\^{}2)/A2} $\to$ real part of $Q_1$\\
        \textbf{D2}: \texttt{=(2*(1+A2)*1 - 2*1*1)/A2} $\to$ imaginary part of $Q_1$
  \item \textbf{Path 2 --- $\Delta z = i\cdot h$ (imaginary):}\\
        $f(z_0+ih) = (1+i(1+h))^2 = 1-(1+h)^2 + 2(1+h)i$\\
        \textbf{E2}: \texttt{=(1-(1+A2)\^{}2 - (1-1\^{}2))/(A2)} $\to$ real part of $Q_2$\\
        \textbf{F2}: \texttt{=(2*(1+A2) - 2*1)/A2} $\to$ imaginary part of $Q_2$
\end{itemize}

\noindent As $h \to 0$, both paths converge to $\operatorname{Re}=2$, $\operatorname{Im}=2$,
confirming $f'(1+i) = 2(1+i) = 2+2i$.

\medskip
\noindent\textbf{Contrast with $g(z) = \overline{z}$:}
For $g$, Path 1 gives quotient $= 1 + 0i$ for all $h$,
while Path 2 gives $-1 + 0i$ for all $h$. These never agree,
confirming non-differentiability regardless of $h$.

\begin{learnbox}{Takeaway --- Excel Example}
Numerical experiments with two orthogonal paths can quickly expose whether a function
is likely differentiable. Convergence to the same value along both paths is a
necessary (but not sufficient) condition for differentiability.
\end{learnbox}

%% ============================================================
\section{Python Example (MANDATORY)}
%% ============================================================

\subsection*{Numerical Comparison: Differentiable vs.\ Non-Differentiable Function}

The following script approximates the difference quotient of $f(z) = z^2$ (differentiable)
and $g(z) = \overline{z}$ (non-differentiable) at $z_0 = 1 + i$ along four paths and
prints the results.

\begin{lstlisting}[language=Python]
import cmath
import numpy as np

z0 = complex(1, 1)          # z0 = 1 + i
h  = 1e-7                   # small step size

def f(z):
    return z**2             # differentiable: f'(z0) = 2*(1+i) = 2+2i

def g(z):
    return z.real - 1j*z.imag   # conjugate: nowhere differentiable

# Four approach directions: real, imaginary, 45-deg, 135-deg
directions = {
    "Real axis (theta=0)":        complex(h, 0),
    "Imag axis (theta=pi/2)":     complex(0, h),
    "45-deg (theta=pi/4)":        complex(h/np.sqrt(2), h/np.sqrt(2)),
    "135-deg (theta=3*pi/4)":     complex(-h/np.sqrt(2), h/np.sqrt(2)),
}

print("=" * 60)
print(f"Difference quotient at z0 = {z0}")
print("=" * 60)

for label, dz in directions.items():
    qf = (f(z0 + dz) - f(z0)) / dz      # quotient for z^2
    qg = (g(z0 + dz) - g(z0)) / dz      # quotient for conj(z)
    print(f"\nDirection: {label}")
    print(f"  f(z)=z^2  => Q = {qf:.6f}   (expected: 2+2i)")
    print(f"  g(z)=zbar => Q = {qg:.6f}   (path-dependent!)")

# Expected output:
# ============================================================
# Difference quotient at z0 = (1+1j)
# ============================================================
#
# Direction: Real axis (theta=0)
#   f(z)=z^2  => Q = (2.000000+2.000000j)   (expected: 2+2i)
#   g(z)=zbar => Q = (1.000000+0.000000j)   (path-dependent!)
#
# Direction: Imag axis (theta=pi/2)
#   f(z)=z^2  => Q = (2.000000+2.000000j)   (expected: 2+2i)
#   g(z)=zbar => Q = (-1.000000+0.000000j)  (path-dependent!)
#
# Direction: 45-deg (theta=pi/4)
#   f(z)=z^2  => Q = (2.000000+2.000000j)   (expected: 2+2i)
#   g(z)=zbar => Q = (0.000000-1.000000j)   (path-dependent!)
#
# Direction: 135-deg (theta=3*pi/4)
#   f(z)=z^2  => Q = (2.000000+2.000000j)   (expected: 2+2i)
#   g(z)=zbar => Q = (0.000000+1.000000j)   (path-dependent!)
\end{lstlisting}

\begin{learnbox}{Takeaway --- Python Example}
For $f(z)=z^2$, the difference quotient $\approx 2+2i$ along every path tested,
confirming differentiability. For $g(z)=\overline{z}$, the quotient depends strongly
on direction, confirming nowhere-differentiability.
\end{learnbox}

%% ============================================================
\section{Viva Questions}
%% ============================================================

\begin{enumerate}[label=\textbf{Q\arabic*.}, leftmargin=2.5em]

\item \textbf{State the definition of the complex derivative of $f$ at $z_0$.}\\
The complex derivative is
$f'(z_0) = \lim_{\Delta z\to 0}\dfrac{f(z_0+\Delta z)-f(z_0)}{\Delta z}$,
provided the limit exists and equals the same value for every path of approach.

\item \textbf{Why is complex differentiability more restrictive than real differentiability?}\\
In $\mathbb{R}$ there are only two directions of approach; the limit must be the same
from left and right. In $\mathbb{C}$ there are infinitely many directions, and the limit
must agree along all of them simultaneously.

\item \textbf{Show that $f(z) = \overline{z}$ is nowhere differentiable.}\\
The difference quotient equals $\overline{\Delta z}/\Delta z$. Along the real axis this
is $1$; along the imaginary axis it is $-1$. Since these differ, the limit does not
exist at any point.

\item \textbf{What is the difference between being differentiable at a point and being analytic at a point?}\\
Differentiable at $z_0$ means the derivative exists at that single point.
Analytic at $z_0$ means differentiable in some open neighbourhood of $z_0$.
Analyticity is the stronger condition.

\item \textbf{State the chain rule for complex functions.}\\
If $g$ is differentiable at $z$ and $f$ is differentiable at $g(z)$, then
$(f \circ g)'(z) = f'(g(z))\cdot g'(z)$.

\item \textbf{What does path independence of the complex derivative imply for $u$ and $v$?}\\
It implies the Cauchy--Riemann equations:
$u_x = v_y$ and $u_y = -v_x$,
where $f = u + iv$ and subscripts denote partial derivatives.

\item \textbf{Is $f(z) = |z|^2$ analytic? Justify.}\\
No. $f$ is differentiable only at $z=0$ (with $f'(0)=0$) and nowhere else.
Since differentiability does not hold on any open set, $f$ is not analytic anywhere.

\item \textbf{State the power rule and give its proof sketch for $n \in \mathbb{Z}$, $n \geq 1$.}\\
The power rule states $\dfrac{d}{dz}(z^n) = nz^{n-1}$.
Proof sketch: expand $(z+\Delta z)^n$ using the binomial theorem;
all terms except $nz^{n-1}\Delta z$ contain $(\Delta z)^k$ for $k\geq 2$,
which vanish after dividing by $\Delta z$ and letting $\Delta z \to 0$.

\end{enumerate}

%% ============================================================
\section{Descriptive Questions}
%% ============================================================

\begin{enumerate}[label=\textbf{D\arabic*.}, leftmargin=2.5em]

\item \textbf{Define the derivative of a complex function at a point. Explain why the path
      of approach of $\Delta z$ is critical, and use two explicit paths to show that
      $f(z)=\overline{z}$ is not differentiable anywhere.} \hfill[10 marks]

\item \textbf{State and prove all standard differentiation rules (sum, product, quotient, chain)
      for complex functions, and illustrate each with a fully worked numerical example.}
      \hfill[10 marks]

\item \textbf{Show from first principles that $\dfrac{d}{dz}(z^n) = nz^{n-1}$ for any positive
      integer $n$. Discuss what happens for negative integers and fractional powers.}
      \hfill[10 marks]

\item \textbf{Prove that differentiability of $f$ at $z_0$ implies continuity of $f$ at $z_0$.
      Show by an example that the converse is false.} \hfill[10 marks]

\item \textbf{Show that the path-independence requirement for the complex derivative naturally
      leads to the Cauchy--Riemann equations $u_x = v_y$ and $u_y = -v_x$.
      Derive these by approaching along the real and imaginary axes.} \hfill[10 marks]

\end{enumerate}

%% ============================================================
\section{Practice Problems}
%% ============================================================

\begin{enumerate}[label=\textbf{P\arabic*.}, leftmargin=2.5em]

\item Find $f'(z)$ from first principles for $f(z) = z^3$.
      \textit{Hint: Expand $(z+\Delta z)^3$ and cancel.}

\item Using the quotient rule, differentiate $h(z) = \dfrac{z^3-1}{z^2+1}$ and
      simplify your answer.
      \textit{Hint: Apply numerator $\cdot$ denominator rule carefully; result involves a
      single fraction.}

\item Show that $f(z) = \operatorname{Im}(z) = y$ is not differentiable at any point by
      computing the difference quotient along two perpendicular paths.
      \textit{Hint: Along the real axis the quotient is $0$; along the imaginary axis it is $1/i = -i$.}

\item Let $f(z) = z|z|$. Determine whether $f$ is differentiable at $z = 0$.
      \textit{Hint: Write $f(\Delta z) = \Delta z \cdot |\Delta z|$; form the quotient
      $\Delta z\,|\Delta z|/\Delta z = |\Delta z| \to 0$, so $f'(0) = 0$.}

\item Differentiate $F(z) = \sin(z^2)$ using the chain rule.
      \textit{Hint: $\dfrac{d}{dz}[\sin(z^2)] = \cos(z^2)\cdot 2z$.}

\item Compute $\dfrac{d}{dz}\!\left[\dfrac{e^z}{z^2+4}\right]$ and find all $z$ where the
      derivative fails to exist.
      \textit{Hint: Use quotient rule; denominator $z^2+4 = 0$ at $z = \pm 2i$.}

\end{enumerate}

%% ============================================================
\section{MCQs}
%% ============================================================

\begin{enumerate}[label=\textbf{MCQ \arabic*.}, leftmargin=3em]

\item The derivative of $f(z)=z^4$ at $z=1+i$ is:
\begin{enumerate}[label=(\alph*)]
  \item $4(1+i)^2$
  \item \textbf{$4(1+i)^3$}
  \item $4(1+i)^4$
  \item $4$
\end{enumerate}
\textit{Explanation:} The power rule gives $f'(z)=4z^3$; substituting $z=1+i$:
$f'(1+i)=4(1+i)^3$.

\medskip
\item Which of the following is \textbf{nowhere differentiable}?
\begin{enumerate}[label=(\alph*)]
  \item $f(z) = z^2$
  \item $f(z) = e^z$
  \item \textbf{$f(z) = \overline{z}$}
  \item $f(z) = \sin z$
\end{enumerate}
\textit{Explanation:} $\overline{z}$ gives path-dependent difference quotients
($1$ vs $-1$); it is nowhere differentiable.

\medskip
\item If the complex derivative $f'(z_0)$ exists, the limit must be independent of:
\begin{enumerate}[label=(\alph*)]
  \item the modulus of $\Delta z$ only
  \item \textbf{the direction (argument) of $\Delta z$}
  \item the real part of $z_0$
  \item the imaginary part of $f(z_0)$
\end{enumerate}
\textit{Explanation:} Path independence means the limit must not depend on the
angle $\arg(\Delta z)$ as $|\Delta z| \to 0$.

\medskip
\item The function $f(z) = |z|^2$ is differentiable:
\begin{enumerate}[label=(\alph*)]
  \item Everywhere in $\mathbb{C}$
  \item Only for $|z| > 1$
  \item \textbf{Only at $z = 0$}
  \item Nowhere
\end{enumerate}
\textit{Explanation:} The difference quotient reduces to a path-dependent expression
at every $z \neq 0$, but at $z=0$ the term $z\cdot(\overline{\Delta z}/\Delta z)$
vanishes and the limit equals $0$.

\medskip
\item Differentiability at a single isolated point implies:
\begin{enumerate}[label=(\alph*)]
  \item Analyticity at that point
  \item \textbf{Continuity at that point}
  \item Analyticity in a neighbourhood
  \item That the Cauchy--Riemann equations hold in an open set
\end{enumerate}
\textit{Explanation:} Differentiability always implies continuity. It does \emph{not}
imply analyticity unless it holds throughout a neighbourhood.

\end{enumerate}

%% ============================================================
\section{Common Mistakes}
%% ============================================================

\begin{mistakebox}{Common Mistakes in Complex Differentiation}
\begin{center}
\begin{tabular}{@{}p{5cm}p{4.5cm}p{5cm}@{}}
\toprule
\textbf{Mistake} & \textbf{Wrong reasoning} & \textbf{Correct reasoning} \\
\midrule
Assuming real differentiability is enough &
``$f(x,y)$ has partial derivatives, so $f$ is complex-differentiable'' &
Complex differentiability requires the Cauchy--Riemann equations AND continuity of partials \\
\midrule
Checking only one path &
``Limit along real axis exists, so derivative exists'' &
Must verify limit is the same along \emph{all} paths; one path is insufficient \\
\midrule
Confusing differentiable with analytic &
``$|z|^2$ is differentiable at $0$, so it's analytic'' &
Analytic requires differentiability in an open neighbourhood, not just one point \\
\midrule
Algebraic error in complex quotient &
Forgetting to multiply numerator and denominator by conjugate when simplifying &
Keep track of real and imaginary parts separately; use $z\overline{z} = |z|^2$ \\
\midrule
Using $\int$ instead of $\oint$ &
Writing $\int_C$ for a closed contour later in the course &
Closed contours always use $\oint_C$; this becomes important in Topics 09--18 \\
\midrule
Ignoring direction of increment &
Computing the quotient only for real $\Delta z$ &
Must set up $\Delta z = \Delta x + i\,\Delta y$ and test arbitrary direction \\
\bottomrule
\end{tabular}
\end{center}
\end{mistakebox}

%% ============================================================
\section{Quick Recap}
%% ============================================================

\begin{learnbox}{Quick Recap --- Topic 03: Differentiation}
\begin{itemize}[leftmargin=1.5em]
  \item The complex derivative is defined by the same difference-quotient formula as in
        real calculus, but the limit must hold along \emph{every} path.
  \item Path independence is the key constraint; it forces
        $u_x = v_y$ and $u_y = -v_x$ (Cauchy--Riemann equations, Topic~04).
  \item All real-calculus rules --- sum, product, quotient, chain, power --- carry over
        verbatim to complex differentiable functions.
  \item $e^z$, $\sin z$, $\cos z$, and all polynomials are entire (differentiable
        everywhere in $\mathbb{C}$).
  \item $\overline{z}$ and $|z|^2$ are \emph{not} analytic; $|z|^2$ is differentiable
        only at the origin.
  \item Differentiability at an isolated point implies continuity, but \emph{not}
        analyticity; analyticity requires differentiability on an open set.
  \item Logarithm branches are analytic only on cut planes; the branch cut must be
        specified and avoided.
  \item The two-path derivation of the Cauchy--Riemann equations is the conceptual
        bridge from this topic to every subsequent result in complex analysis.
\end{itemize}
\end{learnbox}

\end{document}
